\documentclass[a4paper,11pt]{article}
\usepackage[T2A]{fontenc}
\usepackage[utf8]{inputenc}
\usepackage[russian]{babel}
\usepackage{amsmath,amssymb}
\usepackage[margin=23mm]{geometry}
\usepackage{booktabs,tabularx,array}
\usepackage{xcolor,tikz,fancyhdr}
\usepackage[colorlinks=true,linkcolor=teal!70!black]{hyperref}
\definecolor{accent}{HTML}{006C70}
\definecolor{darkaccent}{HTML}{004D4D}
\setlength{\parindent}{0pt}
\setlength{\parskip}{3pt}
\setlength{\headheight}{15pt}
\pagestyle{fancy}
\fancyhf{}
\fancyhead[L]{\small Пропорции}
\fancyhead[R]{\small Надя Клименко}
\fancyfoot[C]{\thepage}
\renewcommand{\headrulewidth}{0.3pt}
\newcommand{\rubric}[1]{\par\smallskip\textbf{\color{darkaccent}#1}\par\nobreak}
\begin{document}
\begin{titlepage}
\centering
\vspace*{25mm}
{\Large\color{darkaccent} МАТЕМАТИКА\par}
\vspace{17mm}
{\LARGE Ревью домашней работы\par}
\vspace{12mm}
{\Large\color{accent} Пропорции\par}
\vfill
{\large Надя Клименко\par}
\vspace{8mm}
28 ноября 2026 г.\par
{\small Домашняя работа от 27 ноября 2026 г.\par}
\vfill
{\small Лёвин Артём Александрович\par эксклюзивно для Нади Клименко\par}
\vspace{10mm}
\begin{tikzpicture}
\draw[accent,line width=.6pt] (0,0) .. controls (2,0.8) and (4,-0.8) .. (6,0) .. controls (8,0.8) and (10,-0.8) .. (12,0);
\end{tikzpicture}
\end{titlepage}
\section*{Сводка замечаний}
\begin{tabularx}{\linewidth}{@{}c p{.20\linewidth} X p{.14\linewidth} p{.14\linewidth}@{}}
\toprule
№ & Тема & Замечание & Ваш ответ & Верный ответ\\
\midrule
\hyperref[task:2]{2} & Смешанные числа, пропорция & Неоднозначная запись числителя при переводе смешанного числа. Последующие вычисления согласованы с верной дробью. & $x=2$ & $x=2$\\
\midrule
\hyperref[task:6]{6} & Проценты от числа & Подмена числа $37{,}5$ на $35{,}5$ в числителе. & $0{,}9$ кг & $0{,}9$ кг\\
\bottomrule
\end{tabularx}
\medskip
В задании № 2 замечание относится к читаемости записи. В задании № 6 обнаружена ошибка при переписывании числа; дальнейшие вычисления выполнены с верным значением.
\par
Проверены все представленные решения.
\clearpage
\thispagestyle{empty}
\vspace*{\fill}
\begin{center}
{\LARGE\color{darkaccent} Разбор замечаний}
\end{center}
\vspace*{\fill}
\clearpage
\phantomsection\label{task:2}
\section*{\color{darkaccent} Задание № 2}
\setlength{\abovedisplayskip}{5pt}
\setlength{\belowdisplayskip}{5pt}
\setlength{\abovedisplayshortskip}{3pt}
\setlength{\belowdisplayshortskip}{3pt}
\rubric{Условие}
Найдите неизвестный член пропорции:
\[2\frac12:1\frac23=3:x.\]
\rubric{Ваше решение и ответ}
Первое смешанное число переведено в дробь $\frac52$. Запись числителя второй дроби в первой строке читается неоднозначно (похожа на 6). В следующей строке уже отчётливо записано $\frac53$:
\[\frac52:\frac53=3:x,\qquad \frac52 x=\frac53\cdot3,\]
\[\frac52 x=5,\qquad x=5:\frac52=5\cdot\frac25=2.\]
Ответ: $x=2$.
\rubric{Где требуется уточнение}
\textcolor{red!75!black}{Замечание:} при переводе числа $1\frac23$ числитель должен быть равен 5. Если записана цифра 6, это локальная ошибка записи: $\frac63=2$, тогда как $1\frac23$ меньше двух. Дальнейший ход решения уже использует правильную дробь. Если в первой строке написана цифра 5, исправлять математическое преобразование не требуется.
\rubric{Ключевая идея}
Чтобы перевести смешанное число в неправильную дробь, умножьте целую часть на знаменатель, прибавьте числитель и сохраните знаменатель:
\[a\frac bc=\frac{ac+b}{c},\qquad c\ne0.\]
\rubric{Исправление и полное решение}
Переведём оба смешанных числа:
\[2\frac12=\frac{2\cdot2+1}{2}=\frac52,\qquad
1\frac23=\frac{1\cdot3+2}{3}=\frac53.\]
По основному свойству пропорции произведение крайних членов равно произведению средних:
\[\frac52 x=\frac53\cdot3=5.\]
Разделим обе части на $\frac52$. При делении на дробь умножаем на обратную ей дробь:
\[x=5:\frac52=5\cdot\frac25=2.\]
\rubric{Проверка}
Число $x=2$ отлично от нуля. Подстановка даёт одинаковые отношения:
\[\frac52:\frac53=\frac52\cdot\frac35=\frac32=3:2.\]
\textcolor{accent}{\textbf{Итог:}} ответ $x=2$ верен. Пишите числители дробей однозначно и последовательно.
\rubric{Задача на закрепление}
Найдите неизвестный член пропорции:
\[3\frac14:1\frac34=7:x.\]
\clearpage
\phantomsection\label{task:6}
\section*{\color{darkaccent} Задание № 6}
\rubric{Условие}
В $2{,}4$ кг сплава содержится $37{,}5\%$ меди. Найдите массу меди.
\rubric{Ваше решение и ответ}
В первой строке записано:
\[37{,}5\%=\frac{35{,}5}{100}=0{,}375.\]
Далее Вы использовали верную десятичную дробь:
\[2{,}4\cdot0{,}375=0{,}9\text{ кг}.\]
Ответ: $0{,}9$ кг.
\rubric{Где возникает ошибка}
\textcolor{red!75!black}{Ошибка:} при переводе процентов в обыкновенную дробь изменилось исходное число. Последнее верно записанное значение — $37{,}5\%$. Первый неверный переход — равенство $37{,}5\%=\frac{35{,}5}{100}$.
\rubric{Почему этот переход неверен}
Один процент — одна сотая целого. Поэтому при переводе процентов в дробь делим исходное число на 100, сохраняя все его цифры. Если разделить $35{,}5$ на 100, получится $0{,}355$. При умножении на массу сплава это дало бы $0{,}852$ кг меди. В Ваших дальнейших вычислениях уже использовано правильное число $0{,}375$.
\rubric{Ключевая идея}
Чтобы найти $p\%$ от числа $A$, переведите процент в долю и умножьте на исходное число:
\[A\cdot\frac{p}{100}.\]
\rubric{Исправление и полное решение}
В данной задаче весь сплав имеет массу $2{,}4$ кг. Доля меди составляет:
\[37{,}5\%=\frac{37{,}5}{100}=0{,}375=\frac38.\]
Умножим массу сплава на эту долю:
\[2{,}4\cdot\frac38=0{,}3\cdot3=0{,}9\text{ кг}.\]
\rubric{Проверка}
Найдём, какой процент от массы сплава составляет полученная масса меди:
\[\frac{0{,}9}{2{,}4}\cdot100\%=\frac38\cdot100\%=37{,}5\%.\]
\textcolor{accent}{\textbf{Итог:}} ответ $0{,}9$ кг верен. Исправьте числитель в первой строке и проверяйте сохранение цифр при переводе процентов в дробь.
\rubric{Задача на закрепление}
В $3{,}2$ кг сплава содержится $42{,}5\%$ меди. Найдите массу меди, записав перевод процентов в десятичную дробь.
\end{document}
